\documentclass[exercice,a4]{thedocument}%
\usepackage[french]{babel}%
\usepackage{pstricks}%
\usepackage[np]{numprint}%
\usepackage{amsmath}%
\startdatakeys
    \source{22GENMATPO3}
    \BO{2026}
    \cycle{4}
    \competencesDuSocle{CA,RA}
    \objectifsApprentissage{C4D2ae,C4A2eb,C4A1aj,C4A2bi,C4B2cc}
    \motsCles{facteurs premiers,fraction,développement,périmètre,équation,volume,pourcentage}
\enddatakeys
\begin{document}%

\medskip

Cet exercice est constitué de six questions indépendantes.

\medskip

\begin{enumerate}
\item Calculer $\dfrac56 + \dfrac78$ et donner le résultat sous la forme d'une fraction irréductible.

On détaillera les calculs.
\item
	\begin{enumerate}
		\item Donner, sans justifier, la décomposition en facteurs premiers de $198$ et de $84$.
		\item En déduire la forme irréductible de la fraction $\dfrac{198}{84}$.
	\end{enumerate}
\item  On donne l'expression littérale suivante : $E = 5(3x - 4) - (2x - 7)$.

Développer et réduire $E$.
\item  On désigne par $b$ un nombre positif.

\begin{minipage}{0.58\linewidth}
Déterminer la valeur de $b$ telle que le périmètre du rectangle ci-contre soit égal à 25.
\end{minipage}\hfill
\begin{minipage}{0.38\linewidth}
\psset{unit=1cm,arrowsize=2pt 3}
\begin{pspicture}(5.8,3.5)
\psframe(1,0)(5.8,2.8)
\psline(1.9,2.7)(1.9,2.9)\psline(2.8,2.7)(2.8,2.9)\psline(3.7,2.7)(3.7,2.9)
\psline(1.4,2.7)(1.5,2.9)\psline(2.3,2.7)(2.4,2.9)\psline(3.2,2.7)(3.3,2.9)
\psline[linewidth=0.6pt]{<->}(1,3)(1.9,3)\uput[u](1.45,3){$b$}
\psline[linewidth=0.6pt]{<->}(3.7,3)(5.8,3)\uput[u](4.75,3){$2,9$}
\psline[linewidth=0.6pt]{<->}(0.8,0)(0.8,2.8)\uput[l](0.8,1.4){4,5}
\end{pspicture}
\end{minipage}

\item ~

\begin{minipage}{0.58\linewidth}
Calculer le volume de la pyramide à base rectangulaire de hauteur SH $= 6$ ci-contre.
\end{minipage}\hfill
\begin{minipage}{0.38\linewidth}
\psset{unit=1cm,arrowsize=2pt 3}
\begin{pspicture}(-2.5,-1)(2.6,4.6)
\psline(0,4.2)(-2.5,-0.5)(0.9,-0.5)(0,4.2)(2.5,0.5)(0.9,-0.5)
\pspolygon[linestyle=dashed](-2.5,-0.5)(2.5,0.5)(-0.9,0.5)
\psline[linestyle=dashed](0,0)(0,4.2)(-0.9,0.5)(0.9,-0.5)
\uput[l](0,2.1){6}\uput[d](-0.8,-0.5){4}\uput[dr](1.7,0){3}
\uput[u](0,4.2){S}\uput[d](0,0){H}
\end{pspicture}
\end{minipage}
\item Le nombre d'habitants d'une ville a augmenté de 12\,\% entre 2019 et 2020.

Cette ville compte \np{20692}~habitants en 2020.

Quel était le nombre d'habitants de cette ville en 2019 ?

\end{enumerate}

\bigskip

\showthesource
\end{document}
